\documentclass{article}

\usepackage[preprint,nonatbib]{project}
\usepackage{enumitem}
\usepackage[utf8]{inputenc}
\usepackage[T1]{fontenc}
\usepackage{url}
\usepackage{booktabs}
\usepackage{amsfonts}
\usepackage{nicefrac}
\usepackage{microtype}
\usepackage{xcolor}
\usepackage{times}
\usepackage{graphicx}
\usepackage{amsmath}
\usepackage{amssymb}
\usepackage{xspace}
\usepackage{tabularx}
\usepackage{multirow}
\usepackage[numbers]{natbib}

\title{How Much Speech Is Enough?\\
Minimal Smartphone Speech Assessment\\
for ALS Dysarthria}

\author{%
  Blake Kemp\\
  Auburn University\\
}

\begin{document}

\maketitle

\begin{abstract}
Amyotrophic lateral sclerosis (ALS) can progressively make speech more difficult as motor control declines. Smartphone voice recordings could provide a practical way to monitor these speech changes remotely, but asking someone with ALS to complete a long set of voice tasks each time may create unnecessary effort. This project looks at whether a smaller set of speech tasks can provide nearly the same performance as the full assessment. I will use the public VOC-ALS dataset, which contains 1,224 recordings from 153 participants completing five sustained vowels and three syllable-repetition tasks. The main analysis focuses on the 102 participants with ALS and uses the ALSFRS-R speech score to define whether speech impairment is present. I will evaluate all 255 possible nonempty combinations of the eight speech tasks using participant-disjoint nested cross-validation. The goal is to find the smallest set of tasks that performs close to the full eight-task assessment while making the assessment easier to complete. This project is focused on remote speech monitoring and is not intended to diagnose ALS.
\end{abstract}

\section{Introduction}

ALS is a progressive neurodegenerative disease that can affect the muscles used for speech. As speech becomes weaker or less clear, repeated clinical or remote assessments may also become harder for someone with ALS to complete. Smartphone recordings are useful because they can be collected outside of a clinic, but a practical system should not ask someone to provide more speech than it actually needs.

This question is personally important to me. My aunt lived with ALS, and her experience previously motivated me to build EyeSpeak-Assist, a low-cost blink-to-speech project. That project focused on helping with communication after normal speech or motor input becomes difficult. This research project looks at an earlier point in that same problem: whether speech changes can be monitored using a shorter and less demanding assessment.

The public VOC-ALS dataset contains eight smartphone speech tasks for each participant: sustained /a/, /e/, /i/, /o/, /u/ and repeated /pa/, /ta/, /ka/ \cite{dubbioso2024vocals}. Rather than only asking which model gives the highest classification performance, this project asks a more practical question: \textit{how many of these tasks are actually necessary?}

The main goal is to identify the smallest subset of tasks that performs nearly as well as the full eight-task protocol when detecting documented speech impairment among participants already diagnosed with ALS. This means balancing two goals: predictive performance and the amount of effort required from the participant. A shorter protocol could be easier to repeat and more practical for future remote monitoring or assistive technology.


\section{Related Work}

The VOC-ALS dataset was introduced by Dubbioso et al. and contains 1,224 smartphone recordings from 153 participants: 102 people with ALS and 51 healthy controls \cite{dubbioso2024vocals}. Each participant completed five sustained vowel tasks and three syllable-repetition tasks. The dataset also includes clinical information and several basic acoustic measurements, including fundamental frequency, jitter, shimmer, and harmonics-to-noise ratio.

Previous studies have shown that speech features can capture useful information about ALS-related speech changes. Simmatis et al. used automatically extracted acoustic features to distinguish people with ALS from healthy controls and to separate different levels of bulbar impairment \cite{simmatis2024acoustic}. More recent work has also compared global and vowel-level acoustic features for dysarthria assessment in motor neuron disease \cite{attia2025global}.

These studies show that speech contains measurable information related to motor impairment. My project focuses on a different question: how much speech is actually needed. Instead of only asking which model performs best when using all available tasks, I will compare every possible subset of the eight VOC-ALS tasks and look at the tradeoff between model performance and the number of tasks required.


\section{Method}

\subsection{Dataset and Prediction Target}

The primary analysis uses the 102 participants with ALS in VOC-ALS. Healthy controls are not included in the main prediction task because the goal of this project is not to diagnose ALS. The target label comes from the ALSFRS-R speech item \cite{cedarbaum1999alsfrsr}. Participants with a speech score of 4 are treated as having normal speech, while scores below 4 are treated as having documented speech impairment.

The eight candidate speech tasks are

$$
T=\{a,e,i,o,u,pa,ta,ka\}.
$$

Since each task can either be included or left out, there are

$$
2^8-1=255
$$

possible nonempty task subsets to evaluate.

\subsection{Acoustic Features}

The first feature set will use the basic acoustic measures already provided with VOC-ALS: mean fundamental frequency, fundamental-frequency variation, harmonics-to-noise ratio, jitter, and shimmer. A second feature set may also be extracted directly from the audio and could include duration, energy, pitch statistics, MFCC summaries, and timing information for /pa/, /ta/, and /ka/. The published features will remain the main baseline so the project stays manageable and easy to reproduce.

\subsection{Models}

Initial model testing will compare regularized logistic regression, support-vector machines, random forests, and gradient-boosted trees. For the full 255-subset experiment, one fixed modeling pipeline will be chosen before looking at the outer-test results. The current leading choice is class-weighted logistic regression because it is simple, easier to interpret, and a good fit for a relatively small dataset.

\subsection{Task Selection}

For each subset $S \subseteq T$, the model will use only the features from the tasks included in $S$. Macro F1 will be the main performance metric, while the number of tasks $|S|$ will represent assessment burden. The goal is to find the smallest subset that performs within $\delta=0.02$ macro F1 of the full eight-task protocol:

$$
Performance(S) \geq Performance(T)-0.02.
$$

I will also test thresholds of 0.01 and 0.05 so the final result is not based on a single cutoff value.

\subsection{Leakage Control and Validation}

All recordings from the same participant will stay in the same fold. The evaluation will use participant-disjoint nested cross-validation. The outer folds will estimate how well the model performs on unseen participants, while the inner folds will be used for model tuning. Scaling, imputation, feature selection, and tuning will only use the training data within each fold. The same outer folds will also be reused for every task subset so the comparisons between subsets are consistent and fair.


\section{Experiment}

The experiment will start by checking the VOC-ALS metadata, confirming that each participant has the expected task recordings, and creating fixed participant-level fold assignments. The current plan is to use five stratified outer folds, with inner participant-level folds for model tuning. If the class sizes make five folds too unstable, I will use a repeated stratified group design instead.

The main experiment will first establish a baseline using all eight speech tasks. Once the model pipeline is fixed, that same pipeline will be evaluated on all 255 nonempty task subsets. Every subset will use the same held-out participants so the results can be compared fairly. The main results will be macro F1 and number of tasks, followed by balanced accuracy, sensitivity, specificity, ROC AUC, PR AUC, calibration, and confusion matrices.

A performance-versus-burden Pareto frontier will be used to show which task combinations give the best tradeoffs between accuracy and assessment length. Participant-level bootstrap confidence intervals and paired bootstrap comparisons will help show whether small differences between task subsets are actually meaningful. I will also track how often each task appears in the best reduced protocols to see whether certain tasks are consistently useful.

At the time of this draft, I am still working on the dataset audit and participant-level evaluation setup, so final model results are not included yet. One possible result is that a two- or three-task protocol performs close to the full eight-task assessment. It is also possible that more tasks are needed, or that there is no single reduced protocol that works best across all folds. Any of these outcomes would still answer the main research question.


\section{Conclusion}

This project asks a simple question with a practical goal: how much speech is actually needed for an ALS speech assessment? Using VOC-ALS, I will compare every possible combination of the eight available smartphone speech tasks while keeping the evaluation participant-disjoint. The goal is not to build an ALS diagnostic tool, but to see whether remote speech monitoring can use a shorter protocol without losing much useful information.

The biggest limitation is the small dataset, especially for participants with more severe speech impairment. Task count is also only a rough estimate of burden and does not directly measure fatigue. Because of these limits, the final report will focus more on uncertainty and patterns that stay consistent across folds than on small score differences. If a shorter task set performs reliably close to the full protocol, it could be a useful starting point for lower-burden smartphone or assistive-technology assessments.



\begin{thebibliography}{9}

\bibitem{dubbioso2024vocals}
R. Dubbioso, M. Spisto, L. Verde, et al.
\newblock Voice signals database of ALS patients with different dysarthria severity and healthy controls.
\newblock \textit{Scientific Data}, 11:800, 2024.

\bibitem{cedarbaum1999alsfrsr}
J. M. Cedarbaum, N. Stambler, E. Malta, et al.
\newblock The ALSFRS-R: a revised ALS functional rating scale that incorporates assessments of respiratory function.
\newblock \textit{Journal of the Neurological Sciences}, 169(1--2):13--21, 1999.

\bibitem{simmatis2024acoustic}
L. E. R. Simmatis, J. Robin, M. J. Spilka, and Y. Yunusova.
\newblock Detecting bulbar amyotrophic lateral sclerosis (ALS) using automatic acoustic analysis.
\newblock \textit{BioMedical Engineering OnLine}, 23:15, 2024.

\bibitem{attia2025global}
M. Attia, L. Simmatis, E. Sejdic, and Y. Yunusova.
\newblock Global vs. segmental acoustic features for dysarthria assessment in motor neuron diseases.
\newblock \textit{Computers in Biology and Medicine}, 199:111291, 2025.

\end{thebibliography}

\newpage
\section*{NeurIPS Paper Checklist}



\begin{enumerate}

\item {\bf Claims}
    \item[] Question: Do the main claims made in the abstract and introduction accurately reflect the paper's contributions and scope?
    \item[] Answer: \answerYes{}
    \item[] Justification: The Abstract and Introduction describe the same goal as the rest of the paper: finding a smaller set of VOC-ALS speech tasks that can still support speech-impairment monitoring. They also make it clear that this project is not trying to diagnose ALS.
    \item[] Guidelines:
    \begin{itemize}
        \item The answer NA means that the abstract and introduction do not include the claims made in the paper.
        \item The abstract and/or introduction should clearly state the claims made, including the contributions made in the paper and important assumptions and limitations. A No or NA answer to this question will not be perceived well by the reviewers. 
        \item The claims made should match theoretical and experimental results, and reflect how much the results can be expected to generalize to other settings. 
        \item It is fine to include aspirational goals as motivation as long as it is clear that these goals are not attained by the paper. 
    \end{itemize}

\item {\bf Limitations}
    \item[] Question: Does the paper discuss the limitations of the work performed by the authors?
    \item[] Answer: \answerYes{}
    \item[] Justification: The Conclusion points out the main limits of the project, including the small dataset, the smaller number of participants with more severe speech impairment, and the fact that task count is only an estimate of effort rather than a direct measure of fatigue.
    \item[] Guidelines:
    \begin{itemize}
        \item The answer NA means that the paper has no limitation while the answer No means that the paper has limitations, but those are not discussed in the paper. 
        \item The authors are encouraged to create a separate "Limitations" section in their paper.
        \item The paper should point out any strong assumptions and how robust the results are to violations of these assumptions (e.g., independence assumptions, noiseless settings, model well-specification, asymptotic approximations only holding locally). The authors should reflect on how these assumptions might be violated in practice and what the implications would be.
        \item The authors should reflect on the scope of the claims made, e.g., if the approach was only tested on a few datasets or with a few runs. In general, empirical results often depend on implicit assumptions, which should be articulated.
        \item The authors should reflect on the factors that influence the performance of the approach. For example, a facial recognition algorithm may perform poorly when image resolution is low or images are taken in low lighting. Or a speech-to-text system might not be used reliably to provide closed captions for online lectures because it fails to handle technical jargon.
        \item The authors should discuss the computational efficiency of the proposed algorithms and how they scale with dataset size.
        \item If applicable, the authors should discuss possible limitations of their approach to address problems of privacy and fairness.
        \item While the authors might fear that complete honesty about limitations might be used by reviewers as grounds for rejection, a worse outcome might be that reviewers discover limitations that aren't acknowledged in the paper. The authors should use their best judgment and recognize that individual actions in favor of transparency play an important role in developing norms that preserve the integrity of the community. Reviewers will be specifically instructed to not penalize honesty concerning limitations.
    \end{itemize}

\item {\bf Theory assumptions and proofs}
    \item[] Question: For each theoretical result, does the paper provide the full set of assumptions and a complete (and correct) proof?
    \item[] Answer: \answerNA{}
    \item[] Justification: This project is focused on an experimental machine-learning study, so there are no new theorems or formal proofs to provide.
    \item[] Guidelines:
    \begin{itemize}
        \item The answer NA means that the paper does not include theoretical results. 
        \item All the theorems, formulas, and proofs in the paper should be numbered and cross-referenced.
        \item All assumptions should be clearly stated or referenced in the statement of any theorems.
        \item The proofs can either appear in the main paper or the supplemental material, but if they appear in the supplemental material, the authors are encouraged to provide a short proof sketch to provide intuition. 
        \item Inversely, any informal proof provided in the core of the paper should be complemented by formal proofs provided in appendix or supplemental material.
        \item Theorems and Lemmas that the proof relies upon should be properly referenced. 
    \end{itemize}

    \item {\bf Experimental result reproducibility}
    \item[] Question: Does the paper fully disclose all the information needed to reproduce the main experimental results of the paper to the extent that it affects the main claims and/or conclusions of the paper (regardless of whether the code and data are provided or not)?
    \item[] Answer: \answerNo{}
    \item[] Justification: The paper explains the overall experiment, including participant-disjoint nested cross-validation, the task-subset search, and the evaluation metrics. However, I have not finalized the exact model pipeline and hyperparameters or run the full experiment yet, so the draft is not fully reproducible yet.
    \item[] Guidelines:
    \begin{itemize}
        \item The answer NA means that the paper does not include experiments.
        \item If the paper includes experiments, a No answer to this question will not be perceived well by the reviewers: Making the paper reproducible is important, regardless of whether the code and data are provided or not.
        \item If the contribution is a dataset and/or model, the authors should describe the steps taken to make their results reproducible or verifiable. 
        \item Depending on the contribution, reproducibility can be accomplished in various ways. For example, if the contribution is a novel architecture, describing the architecture fully might suffice, or if the contribution is a specific model and empirical evaluation, it may be necessary to either make it possible for others to replicate the model with the same dataset, or provide access to the model. In general. releasing code and data is often one good way to accomplish this, but reproducibility can also be provided via detailed instructions for how to replicate the results, access to a hosted model (e.g., in the case of a large language model), releasing of a model checkpoint, or other means that are appropriate to the research performed.
        \item While NeurIPS does not require releasing code, the conference does require all submissions to provide some reasonable avenue for reproducibility, which may depend on the nature of the contribution. For example
        \begin{enumerate}
            \item If the contribution is primarily a new algorithm, the paper should make it clear how to reproduce that algorithm.
            \item If the contribution is primarily a new model architecture, the paper should describe the architecture clearly and fully.
            \item If the contribution is a new model (e.g., a large language model), then there should either be a way to access this model for reproducing the results or a way to reproduce the model (e.g., with an open-source dataset or instructions for how to construct the dataset).
            \item We recognize that reproducibility may be tricky in some cases, in which case authors are welcome to describe the particular way they provide for reproducibility. In the case of closed-source models, it may be that access to the model is limited in some way (e.g., to registered users), but it should be possible for other researchers to have some path to reproducing or verifying the results.
        \end{enumerate}
    \end{itemize}


\item {\bf Open access to data and code}
    \item[] Question: Does the paper provide open access to the data and code, with sufficient instructions to faithfully reproduce the main experimental results, as described in supplemental material?
    \item[] Answer: \answerNo{}
    \item[] Justification: VOC-ALS is a public dataset and is cited in the paper, but I have not released the project code and reproduction scripts with this draft.
    \item[] Guidelines:
    \begin{itemize}
        \item The answer NA means that paper does not include experiments requiring code.
        \item Please see the NeurIPS code and data submission guidelines (\url{https://nips.cc/public/guides/CodeSubmissionPolicy}) for more details.
        \item While we encourage the release of code and data, we understand that this might not be possible, so “No” is an acceptable answer. Papers cannot be rejected simply for not including code, unless this is central to the contribution (e.g., for a new open-source benchmark).
        \item The instructions should contain the exact command and environment needed to run to reproduce the results. See the NeurIPS code and data submission guidelines (\url{https://nips.cc/public/guides/CodeSubmissionPolicy}) for more details.
        \item The authors should provide instructions on data access and preparation, including how to access the raw data, preprocessed data, intermediate data, and generated data, etc.
        \item The authors should provide scripts to reproduce all experimental results for the new proposed method and baselines. If only a subset of experiments are reproducible, they should state which ones are omitted from the script and why.
        \item At submission time, to preserve anonymity, the authors should release anonymized versions (if applicable).
        \item Providing as much information as possible in supplemental material (appended to the paper) is recommended, but including URLs to data and code is permitted.
    \end{itemize}


\item {\bf Experimental setting/details}
    \item[] Question: Does the paper specify all the training and test details (e.g., data splits, hyperparameters, how they were chosen, type of optimizer, etc.) necessary to understand the results?
    \item[] Answer: \answerNo{}
    \item[] Justification: The Method and Experiment sections explain the planned data splits, candidate models, metrics, and all 255 task subsets. The final hyperparameters and fixed pipeline are still being finalized before the full experiment is run.
    \item[] Guidelines:
    \begin{itemize}
        \item The answer NA means that the paper does not include experiments.
        \item The experimental setting should be presented in the core of the paper to a level of detail that is necessary to appreciate the results and make sense of them.
        \item The full details can be provided either with the code, in appendix, or as supplemental material.
    \end{itemize}

\item {\bf Experiment statistical significance}
    \item[] Question: Does the paper report error bars suitably and correctly defined or other appropriate information about the statistical significance of the experiments?
    \item[] Answer: \answerNo{}
    \item[] Justification: I have not run the full experiment yet, so there are no final error bars or significance results to report. The Experiment section states that I plan to use participant-level bootstrap confidence intervals and paired bootstrap comparisons.
    \item[] Guidelines:
    \begin{itemize}
        \item The answer NA means that the paper does not include experiments.
        \item The authors should answer "Yes" if the results are accompanied by error bars, confidence intervals, or statistical significance tests, at least for the experiments that support the main claims of the paper.
        \item The factors of variability that the error bars are capturing should be clearly stated (for example, train/test split, initialization, random drawing of some parameter, or overall run with given experimental conditions).
        \item The method for calculating the error bars should be explained (closed form formula, call to a library function, bootstrap, etc.)
        \item The assumptions made should be given (e.g., Normally distributed errors).
        \item It should be clear whether the error bar is the standard deviation or the standard error of the mean.
        \item It is OK to report 1-sigma error bars, but one should state it. The authors should preferably report a 2-sigma error bar than state that they have a 96\% CI, if the hypothesis of Normality of errors is not verified.
        \item For asymmetric distributions, the authors should be careful not to show in tables or figures symmetric error bars that would yield results that are out of range (e.g. negative error rates).
        \item If error bars are reported in tables or plots, The authors should explain in the text how they were calculated and reference the corresponding figures or tables in the text.
    \end{itemize}

\item {\bf Experiments compute resources}
    \item[] Question: For each experiment, does the paper provide sufficient information on the computer resources (type of compute workers, memory, time of execution) needed to reproduce the experiments?
    \item[] Answer: \answerNo{}
    \item[] Justification: I have not documented the final hardware, memory, run times, or total compute yet because the full experiment has not been run.
    \item[] Guidelines:
    \begin{itemize}
        \item The answer NA means that the paper does not include experiments.
        \item The paper should indicate the type of compute workers CPU or GPU, internal cluster, or cloud provider, including relevant memory and storage.
        \item The paper should provide the amount of compute required for each of the individual experimental runs as well as estimate the total compute. 
        \item The paper should disclose whether the full research project required more compute than the experiments reported in the paper (e.g., preliminary or failed experiments that didn't make it into the paper). 
    \end{itemize}
    
\item {\bf Code of ethics}
    \item[] Question: Does the research conducted in the paper conform, in every respect, with the NeurIPS Code of Ethics \url{https://neurips.cc/public/EthicsGuidelines}?
    \item[] Answer: \answerNA{}
    \item[] Justification: I have not done a separate review of the NeurIPS Code of Ethics for this course draft. The project is limited to assistive speech monitoring and does not claim to diagnose ALS or replace clinical care.
    \item[] Guidelines:
    \begin{itemize}
        \item The answer NA means that the authors have not reviewed the NeurIPS Code of Ethics.
        \item If the authors answer No, they should explain the special circumstances that require a deviation from the Code of Ethics.
        \item The authors should make sure to preserve anonymity (e.g., if there is a special consideration due to laws or regulations in their jurisdiction).
    \end{itemize}


\item {\bf Broader impacts}
    \item[] Question: Does the paper discuss both potential positive societal impacts and negative societal impacts of the work performed?
    \item[] Answer: \answerNo{}
    \item[] Justification: The paper explains the possible benefit of a shorter, lower-burden speech assessment, but I have not yet added a full discussion of possible downsides such as privacy, bias, or misuse.
    \item[] Guidelines:
    \begin{itemize}
        \item The answer NA means that there is no societal impact of the work performed.
        \item If the authors answer NA or No, they should explain why their work has no societal impact or why the paper does not address societal impact.
        \item Examples of negative societal impacts include potential malicious or unintended uses (e.g., disinformation, generating fake profiles, surveillance), fairness considerations (e.g., deployment of technologies that could make decisions that unfairly impact specific groups), privacy considerations, and security considerations.
        \item The conference expects that many papers will be foundational research and not tied to particular applications, let alone deployments. However, if there is a direct path to any negative applications, the authors should point it out. For example, it is legitimate to point out that an improvement in the quality of generative models could be used to generate deepfakes for disinformation. On the other hand, it is not needed to point out that a generic algorithm for optimizing neural networks could enable people to train models that generate Deepfakes faster.
        \item The authors should consider possible harms that could arise when the technology is being used as intended and functioning correctly, harms that could arise when the technology is being used as intended but gives incorrect results, and harms following from (intentional or unintentional) misuse of the technology.
        \item If there are negative societal impacts, the authors could also discuss possible mitigation strategies (e.g., gated release of models, providing defenses in addition to attacks, mechanisms for monitoring misuse, mechanisms to monitor how a system learns from feedback over time, improving the efficiency and accessibility of ML).
    \end{itemize}
    
\item {\bf Safeguards}
    \item[] Question: Does the paper describe safeguards that have been put in place for responsible release of data or models that have a high risk for misuse (e.g., pretrained language models, image generators, or scraped datasets)?
    \item[] Answer: \answerNA{}
    \item[] Justification: This project is not releasing a high-risk model or scraped dataset, so special release safeguards do not apply.
    \item[] Guidelines:
    \begin{itemize}
        \item The answer NA means that the paper poses no such risks.
        \item Released models that have a high risk for misuse or dual-use should be released with necessary safeguards to allow for controlled use of the model, for example by requiring that users adhere to usage guidelines or restrictions to access the model or implementing safety filters. 
        \item Datasets that have been scraped from the Internet could pose safety risks. The authors should describe how they avoided releasing unsafe images.
        \item We recognize that providing effective safeguards is challenging, and many papers do not require this, but we encourage authors to take this into account and make a best faith effort.
    \end{itemize}

\item {\bf Licenses for existing assets}
    \item[] Question: Are the creators or original owners of assets (e.g., code, data, models), used in the paper, properly credited and are the license and terms of use explicitly mentioned and properly respected?
    \item[] Answer: \answerNo{}
    \item[] Justification: I cite the original VOC-ALS dataset, but I have not yet added its exact license and terms of use to the paper.
    \item[] Guidelines:
    \begin{itemize}
        \item The answer NA means that the paper does not use existing assets.
        \item The authors should cite the original paper that produced the code package or dataset.
        \item The authors should state which version of the asset is used and, if possible, include a URL.
        \item The name of the license (e.g., CC-BY 4.0) should be included for each asset.
        \item For scraped data from a particular source (e.g., website), the copyright and terms of service of that source should be provided.
        \item If assets are released, the license, copyright information, and terms of use in the package should be provided. For popular datasets, \url{paperswithcode.com/datasets} has curated licenses for some datasets. Their licensing guide can help determine the license of a dataset.
        \item For existing datasets that are re-packaged, both the original license and the license of the derived asset (if it has changed) should be provided.
        \item If this information is not available online, the authors are encouraged to reach out to the asset's creators.
    \end{itemize}

\item {\bf New assets}
    \item[] Question: Are new assets introduced in the paper well documented and is the documentation provided alongside the assets?
    \item[] Answer: \answerNA{}
    \item[] Justification: This draft does not release a new dataset, trained model, or software package, so there are no new research assets to document yet.
    \item[] Guidelines:
    \begin{itemize}
        \item The answer NA means that the paper does not release new assets.
        \item Researchers should communicate the details of the dataset/code/model as part of their submissions via structured templates. This includes details about training, license, limitations, etc. 
        \item The paper should discuss whether and how consent was obtained from people whose asset is used.
        \item At submission time, remember to anonymize your assets (if applicable). You can either create an anonymized URL or include an anonymized zip file.
    \end{itemize}

\item {\bf Crowdsourcing and research with human subjects}
    \item[] Question: For crowdsourcing experiments and research with human subjects, does the paper include the full text of instructions given to participants and screenshots, if applicable, as well as details about compensation (if any)? 
    \item[] Answer: \answerNA{}
    \item[] Justification: I am not recruiting or interacting with participants. The project uses the existing public VOC-ALS dataset, so there are no new participant instructions or compensation for this study.
    \item[] Guidelines:
    \begin{itemize}
        \item The answer NA means that the paper does not involve crowdsourcing nor research with human subjects.
        \item Including this information in the supplemental material is fine, but if the main contribution of the paper involves human subjects, then as much detail as possible should be included in the main paper. 
        \item According to the NeurIPS Code of Ethics, workers involved in data collection, curation, or other labor should be paid at least the minimum wage in the country of the data collector. 
    \end{itemize}

\item {\bf Institutional review board (IRB) approvals or equivalent for research with human subjects}
    \item[] Question: Does the paper describe potential risks incurred by study participants, whether such risks were disclosed to the subjects, and whether Institutional Review Board (IRB) approvals (or an equivalent approval/review based on the requirements of your country or institution) were obtained?
    \item[] Answer: \answerNA{}
    \item[] Justification: I am not collecting new human-subject data or recruiting participants; this is secondary analysis of an existing public dataset. I still need to confirm whether an exemption or other formal determination is required before the final report.
    \item[] Guidelines:
    \begin{itemize}
        \item The answer NA means that the paper does not involve crowdsourcing nor research with human subjects.
        \item Depending on the country in which research is conducted, IRB approval (or equivalent) may be required for any human subjects research. If you obtained IRB approval, you should clearly state this in the paper. 
        \item We recognize that the procedures for this may vary significantly between institutions and locations, and we expect authors to adhere to the NeurIPS Code of Ethics and the guidelines for their institution. 
        \item For initial submissions, do not include any information that would break anonymity (if applicable), such as the institution conducting the review.
    \end{itemize}

\item {\bf Declaration of LLM usage}
    \item[] Question: Does the paper describe the usage of LLMs if it is an important, original, or non-standard component of the core methods in this research? Note that if the LLM is used only for writing, editing, or formatting purposes and does not impact the core methodology, scientific rigorousness, or originality of the research, declaration is not required.
    %this research? 
    \item[] Answer: \answerNA{}
    \item[] Justification: LLMs are not part of the feature extraction, prediction models, task selection, or evaluation. Any LLM use is limited to writing or formatting and does not change the research method.
    \item[] Guidelines:
    \begin{itemize}
        \item The answer NA means that the core method development in this research does not involve LLMs as any important, original, or non-standard components.
        \item Please refer to our LLM policy (\url{https://neurips.cc/Conferences/2025/LLM}) for what should or should not be described.
    \end{itemize}

\end{enumerate}


\end{document}
